Este trabajo analiza el conjunto de datos Boston del paquete MASS siguiendo las fases 1 a 5 de CRISP-DM y usando An\'alisis de Componentes Principales como t\'ecnica central. Primero se describen las variables, se verifican valores faltantes y se detectan valores at\'ipicos en varias columnas, con presencia clara en 
\textit{black}, 
\textit{zn}, 
\textit{crim}, 
\textit{medv}, 
\textit{chas}, 
\textit{rm}, 
\textit{ptratio}, 
\textit{lstat} y 
\textit{dis}. Luego se estandarizan las variables explicativas y se ajusta un PCA, donde el primer componente explica 51.06\% de la varianza y cinco componentes alcanzan 80\% acumulado. En la fase de evaluaci\'on se responde que el tama\~no de la vivienda influye positivamente en el precio, que la cercan\'ia al r\'io Charles tambi\'en se asocia con valores mayores, que el estatus socioecon\'omico tiene efecto negativo marcado y que la criminalidad puede modelarse con variables socioecon\'omicas y espaciales. El informe integra evidencia gr\'afica y num\'erica para sustentar cada conclusi\'on.